\section{سری لورنت، نقاط تکین و مانده‌ها}
\def\EMfootchapter{فصل ۱۳ \ \textbullet\ سری لورنت، نقاط تکین و مانده‌ها}
\def\EMtitlesize{24}
\begin{frame}{فصل ۱۳: سری لورنت، نقاط تکین و مانده‌ها}
{\fontsize{20}{26}\selectfont\bfseries\color{SLdef}\EMalign مطالب این فصل\par}
\vspace{10bp}
\EMitem{\textbullet}{\dimexpr\linewidth-26bp\relax}{سری لورنت}
\EMitem{\textbullet}{\dimexpr\linewidth-26bp\relax}{مانده}
\EMitem{\textbullet}{\dimexpr\linewidth-26bp\relax}{قضیهٔ مانده‌ها}
\EMitem{\textbullet}{\dimexpr\linewidth-26bp\relax}{مثال‌هایی از بسط تیلور و لورنت}
\EMitem{\textbullet}{\dimexpr\linewidth-26bp\relax}{مثال‌هایی از نقاط تکین و مانده}
\end{frame}

\def\EMtitlesize{24}
\begin{frame}{سری لورنت}
\begin{EMbox}{thm}{1}{}
\begin{uncoverenv}<2->
\EMp{اگر تابع \EMim{f(z)} تابعی تحلیلی در ناحیهٔ \EMim{D} باشد، سری لورنت این تابع حول نقطهٔ \EMim{z_0} به صورت زیر نوشته می‌شود:}
\vspace{3bp}
\EMeq{f(z) =\sum_{n=-\infty}^{\infty}a_n\,(z-z_0)^n,\qquad a_n =\frac1{2\pi i}\oint_C\frac{f(v)}{(v-z_0)^{n+1}}\,dv}
\end{uncoverenv}
\EMgap
\begin{uncoverenv}<3->
\EMp{که در آن \EMim{C} دایره‌ای به مرکز \EMim{z_0} است و تابع در همهٔ نقاط داخل و روی این دایره، مگر استثنائاً در \EMim{z_0} تحلیلی است.}
\end{uncoverenv}
\EMgap
\begin{uncoverenv}<4->
\EMp{شعاع همگرائی، فاصلهٔ \EMim{z_0} تا نزدیک‌ترین نقطهٔ غیرتحلیلی تابع \EMim{f(z)} است.}
\end{uncoverenv}
\end{EMbox}
\begin{uncoverenv}<5->
\EMp{برای \EMim{n} های نامنفی داریم:}
\vspace{3bp}
\EMeq{a_n =\frac1{2\pi i}\oint_C\frac{f(v)}{(v-z_0)^{n+1}}\,dv =\frac{f^{(n)}(z_0)}{n!},\qquad n =0,1,2,\dots}
\end{uncoverenv}
\end{frame}

\def\EMtitlesize{24}
\begin{frame}{سری لورنت (ادامه)}
\begin{uncoverenv}<2->
\EMp{بنابراین خواهیم داشت:}
\vspace{3bp}
\EMeq{\begin{aligned}
f(z)&=\sum_{n=0}^{\infty}\frac{f^{(n)}(z_0)}{n!}(z-z_0)^n+\sum_{n=-\infty}^{-1}a_n(z-z_0)^n\\
&=\sum_{n=0}^{\infty}\frac{f^{(n)}(z_0)}{n!}(z-z_0)^n+\sum_{n=1}^{\infty}a_{-n}(z-z_0)^{-n}\\
&=\sum_{n=0}^{\infty}\frac{f^{(n)}(z_0)}{n!}(z-z_0)^n+\sum_{n=1}^{\infty}\frac{A_n}{(z-z_0)^n},\qquad A_n=a_{-n}
\end{aligned}}
\end{uncoverenv}
\EMgap
\begin{uncoverenv}<3->
\EMeq{A_n =a_{-n} =\frac1{2\pi i}\oint_C\frac{f(v)}{(v-z_0)^{-n+1}}\,dv =\frac1{2\pi i}\oint_C(v-z_0)^{n-1}f(v)\,dv,\qquad n =1,2,\dots}
\end{uncoverenv}
\end{frame}

\def\EMtitlesize{24}
\begin{frame}{سری لورنت (ادامه)}
\begin{EMbox}{imp}{1}{}
\begin{uncoverenv}<2->
\EMeq{f(z) =\underbrace{\sum_{n=0}^{\infty}\frac{f^{(n)}(z_0)}{n!}(z-z_0)^n}_{\EMt{بخش سری تیلور}}+\underbrace{\sum_{n=1}^{\infty}\frac{A_n}{(z-z_0)^n}}_{\EMt{بخش تکین}},\qquad A_n =\frac1{2\pi i}\oint_C(v-z_0)^{n-1}f(v)\,dv}
\end{uncoverenv}
\end{EMbox}
\begin{uncoverenv}<3->
\EMp{در صورتی که \EMim{z_0} نقطه‌ای غیر عادی (تکین) باشد، به عبارت دیگر تابع در این نقطه تحلیلی نباشد، بخش سری تیلور برای بیان تابع به تنهائی کافی نیست.}
\end{uncoverenv}
\end{frame}

\def\EMtitlesize{24}
\begin{frame}{حالت‌های مختلف نقطهٔ \EMim{z_0}}
\begin{uncoverenv}<2->
\EMitem{۱.}{\dimexpr\linewidth-26bp\relax}{اگر تمامی \EMim{A_n} ها صفر باشند، نقطهٔ \EMim{z_0} یک نقطهٔ \textbf{معمولی} (\lr{Regular}) نامیده می‌شود. سری لورنت در این حالت همان سری تیلور خواهد بود.}
\end{uncoverenv}
\EMgap
\begin{uncoverenv}<3->
\EMitem{۲.}{\dimexpr\linewidth-26bp\relax}{اگر \EMim{A_p\neq0} ولی \EMim{(A_n=0\ \ \forall n>p)}، آنگاه \EMim{z_0} را \textbf{قطب} (\lr{pole}) \textbf{از مرتبهٔ \EMim{p}} می‌نامند.}
\end{uncoverenv}
\EMgap
\begin{uncoverenv}<4->
\EMitem{۳.}{\dimexpr\linewidth-26bp\relax}{اگر \EMim{z_0} نقطهٔ عادی نبوده و ضمناً قطب نیز نباشد، (به عبارت دیگر مقداری مانند \EMim{p} یافت نشود به قسمی که \EMim{A_p\neq0} ولی \EMim{(A_n=0\ \ \forall n>p)})، آنگاه \EMim{z_0} را نقطهٔ \textbf{تکین اصلی} (\lr{Essential Singularity}) می‌نامند.}
\end{uncoverenv}
\end{frame}

\def\EMtitlesize{24}
\begin{frame}{مانده}
\begin{EMbox}{def}{1}{}
\begin{uncoverenv}<2->
\EMp{مانده تابع \EMim{f(z)} در نقطهٔ \EMim{z_0}: \EMim{A_1}.}
\vspace{3bp}
\EMeq{\operatorname{Res}\{f(z)\}\Big|_{z=z_0} =A_1 =\frac1{2\pi i}\oint_Cf(v)\,dv}
\end{uncoverenv}
\end{EMbox}
\begin{EMbox}{thm}{2}{}
\begin{uncoverenv}<3->
\EMp{اگر تابع \EMim{f(z)} در نقطهٔ \EMim{z_0} دارای قطب از مرتبهٔ \EMim{p} باشد داریم:}
\vspace{3bp}
\EMeq{A_k =\frac1{(p-k)!}\left.\frac{d^{\,p-k}}{dz^{\,p-k}}\Big[(z-z_0)^p\,f(z)\Big]\right|_{z=z_0},\qquad k =1,2,\dots,p}
\end{uncoverenv}
\end{EMbox}
\begin{EMbox}{prf}{3}{}
\begin{uncoverenv}<4->
\EMeq{f(z) =\sum_{n=1}^{\infty}\frac{A_n}{(z-z_0)^n}+\sum_{n=0}^{\infty}\frac{f^{(n)}(z_0)}{n!}(z-z_0)^n}
\end{uncoverenv}
\EMgap
\begin{uncoverenv}<5->
\EMp{اگر تابع \EMim{f(z)} در نقطهٔ \EMim{z_0} دارای قطب از مرتبهٔ \EMim{p} باشد داریم:}
\vspace{3bp}
\EMeq{f(z) =\frac{A_p}{(z-z_0)^p}+\frac{A_{p-1}}{(z-z_0)^{p-1}}+\cdots+\frac{A_2}{(z-z_0)^2}+\frac{A_1}{z-z_0}+\sum_{n=0}^{\infty}\frac{f^{(n)}(z_0)}{n!}(z-z_0)^n}
\end{uncoverenv}
\end{EMbox}
\end{frame}

\def\EMtitlesize{24}
\begin{frame}{مانده (ادامه)}
\begin{EMbox}{prf}{1}{}
\begin{uncoverenv}<2->
\EMp{اگر طرفین را در \EMim{(z-z_0)^p} ضرب کنیم داریم:}
\vspace{3bp}
\EMeq{(z-z_0)^pf(z) =A_p+A_{p-1}(z-z_0)+\cdots+A_2(z-z_0)^{p-2}+A_1(z-z_0)^{p-1}+\sum_{n=0}^{\infty}\frac{f^{(n)}(z_0)}{n!}(z-z_0)^{n+p}}
\end{uncoverenv}
\EMgap
\begin{uncoverenv}<3->
\EMp{حال اگر به \EMim{z} مقدار \EMim{z_0} بدهیم، داریم:}
\vspace{3bp}
\EMeq{(z-z_0)^pf(z)\Big|_{z=z_0} =A_p}
\end{uncoverenv}
\EMgap
\begin{uncoverenv}<4->
\EMp{اگر از طرفین \EMim{(z-z_0)^pf(z)} یک بار مشتق بگیریم و به \EMim{z} مقدار \EMim{z_0} بدهیم، داریم:}
\vspace{3bp}
\EMeq{\begin{aligned}
\frac{d}{dz}\Big[(z-z_0)^pf(z)\Big]\Big|_{z=z_0}&=\Bigg\{A_{p-1}+\cdots+A_2(p-2)(z-z_0)^{p-3}+A_1(p-1)(z-z_0)^{p-2}\\
&\qquad+\sum_{n=0}^{\infty}\frac{f^{(n)}(z_0)}{n!}(z-z_0)^{n+p-1}(n+p)\Bigg\}\Bigg|_{z=z_0}=A_{p-1}
\end{aligned}}
\end{uncoverenv}
\end{EMbox}
\end{frame}

\def\EMtitlesize{24}
\begin{frame}{مانده (ادامه)}
\begin{EMbox}{prf}{1}{}
\begin{uncoverenv}<2->
\EMp{اگر از طرفین \EMim{(z-z_0)^pf(z)} دو بار مشتق بگیریم و به \EMim{z} مقدار \EMim{z_0} بدهیم، داریم:}
\vspace{3bp}
\EMeq{\begin{aligned}
\frac{d^2}{dz^2}\Big[(z-z_0)^pf(z)\Big]\Big|_{z=z_0}&=\Bigg\{2A_{p-2}+\cdots+A_2(p-2)(p-3)(z-z_0)^{p-4}+A_1(p-1)(p-2)(z-z_0)^{p-3}\\
&\qquad+\sum_{n=0}^{\infty}\frac{f^{(n)}(z_0)}{n!}(z-z_0)^{n+p-2}(n+p)(n+p-1)\Bigg\}\Bigg|_{z=z_0}=2A_{p-2}
\end{aligned}}
\end{uncoverenv}
\EMgap
\begin{uncoverenv}<3->
\EMeq{\Longrightarrow\quad A_{p-2} =\frac1{2!}\left.\frac{d^2}{dz^2}\Big[(z-z_0)^pf(z)\Big]\right|_{z=z_0}}
\end{uncoverenv}
\end{EMbox}
\end{frame}

\def\EMtitlesize{24}
\begin{frame}{مانده (ادامه)}
\begin{EMbox}{exe}{1}{}
\EMp{به همین ترتیب مشتق‌گیری را ادامه دهید و اثبات قضیه را کامل کنید.}
\end{EMbox}
\begin{EMbox}{rem}{1}{}
\begin{uncoverenv}<2->
\EMp{در حالت خاص، در صورتی که تابع \EMim{f(z)} در نقطهٔ \EMim{z_0} دارای قطب از مرتبهٔ \EMim{p} باشد، مقدار مانده از رابطهٔ زیر به دست می‌آید:}
\vspace{3bp}
\EMeq{\operatorname{Res}\{f(z)\} =A_1 =\frac1{(p-1)!}\left.\frac{d^{\,p-1}}{dz^{\,p-1}}\Big[(z-z_0)^p\,f(z)\Big]\right|_{z=z_0}}
\end{uncoverenv}
\end{EMbox}
\end{frame}

\def\EMtitlesize{24}
\begin{frame}{قضیهٔ مانده‌ها}
\begin{minipage}[t]{500bp}\vspace{0pt}
\begin{EMbox}{thm}{1}{}
\begin{uncoverenv}<2->
\EMp{اگر تابع \EMim{f(z)} در همه‌جای داخل مسیر \EMim{C} به‌جز نقاط \EMim{z_1,z_2,\dots,z_n} تحلیلی باشد و مانده تابع در این نقاط به ترتیب \EMim{A_1^1,A_1^2,\dots,A_1^n} باشد، آنگاه:}
\vspace{3bp}
\EMeq{\oint_Cf(z)\,dz =2\pi i\sum_{k=1}^{n}A_1^k}
\end{uncoverenv}
\end{EMbox}
\EMgap
\begin{uncoverenv}<4->
\EMp{اثبات به عنوان تمرین بر عهدهٔ خواننده است.}
\end{uncoverenv}
\end{minipage}
\EMside{157.4bp}{\uncover<3->{\EMfigure{1.500}{m13-residue-theorem}{مسیر \EMim{C} و نقاط تکین \EMim{z_1,\dots,z_n} در داخل آن؛ حول هر نقطه دایره‌ای کوچک رسم شده است}}}
\end{frame}

\def\EMtitlesize{24}
\begin{frame}{مثال‌هایی از بسط تیلور و لورنت}
\begin{EMbox}{ex}{1}{}
\EMp{مطلوب است محاسبهٔ سری لورنت تابع \EMim{f(z)=\dfrac1{z^4+1}} در همسایگی نقطهٔ \EMim{z_0=0}.}
\EMgap
\begin{uncoverenv}<2->
\EMp{دیده می‌شود که تابع \EMim{f(z)} در نقطهٔ \EMim{z_0=0} تحلیلی است و لذا این نقطه یک نقطهٔ معمولی است و سری لورنت همان سری تیلور (مک‌لورن) خواهد بود. لذا داریم:}
\vspace{3bp}
\EMeq{f(z) =f(0)+\frac1{1!}f^{(1)}(0)z+\frac1{2!}f^{(2)}(0)z^2+\frac1{3!}f^{(3)}(0)z^3+\cdots}
\end{uncoverenv}
\EMgap
\begin{uncoverenv}<3->
\EMp{سری هندسی:}
\vspace{3bp}
\EMeq{1+q+q^2+q^3+q^4+\cdots =\frac1{1-q},\qquad |q|<1}
\end{uncoverenv}
\EMgap
\begin{uncoverenv}<4->
\EMeq{\frac1{1+z^4} =1-z^4+z^8-z^{12}+-\cdots,\qquad |z^4|<1\  \Longrightarrow\ |z|<1}
\end{uncoverenv}
\EMgap
\begin{uncoverenv}<5->
\EMp{(سری تیلور/مک‌لورن؛ ناحیهٔ همگرائی با شعاع ۱.)}
\vspace{3bp}
\EMeq{z^4+1 =0\; \Longrightarrow\;z_1 =e^{i\pi/4},\quad z_2 =e^{i3\pi/4},\quad z_3 =e^{i5\pi/4},\quad z_4 =e^{i7\pi/4}}
\end{uncoverenv}
\EMgap
\begin{uncoverenv}<6->
\EMp{مشاهده می‌شود که فاصلهٔ \EMim{z_0=0} تا نزدیک‌ترین نقطهٔ غیرتحلیلی تابع (شعاع همگرائی) عدد ۱ می‌باشد.}
\end{uncoverenv}
\end{EMbox}
\end{frame}

\def\EMtitlesize{24}
\begin{frame}{مثال‌هایی از بسط تیلور و لورنت (ادامه)}
\begin{minipage}[t]{500bp}\vspace{0pt}
\end{minipage}
\EMside{0.0bp}{\uncover<2->{\EMfigure{1.500}{m13-radius-1}{چهار قطب \EMim{z^4+1=0} روی دایرهٔ واحد؛ شعاع همگرائی سری حول مبدأ برابر ۱ است}}}
\end{frame}

\def\EMtitlesize{24}
\begin{frame}{مثال‌هایی از بسط تیلور و لورنت (ادامه)}
\begin{EMbox}{ex}{1}{}
\EMp{مطلوب است محاسبهٔ سری لورنت تابع \EMim{f(z)=\dfrac{4z^2+30z+68}{(z+4)^2(z-2)}} در همسایگی نقطهٔ \EMim{z_0=0}.}
\EMgap
\begin{uncoverenv}<2->
\EMp{دیده می‌شود که تابع \EMim{f(z)} در نقطهٔ \EMim{z_0=0} تحلیلی است و لذا این نقطه یک نقطهٔ معمولی است و سری لورنت همان سری تیلور (مک‌لورن) خواهد بود. لذا داریم:}
\vspace{3bp}
\EMeq{f(z) =f(0)+\frac1{1!}f^{(1)}(0)z+\frac1{2!}f^{(2)}(0)z^2+\frac1{3!}f^{(3)}(0)z^3+\cdots}
\end{uncoverenv}
\EMgap
\begin{uncoverenv}<3->
\EMeq{\frac1{z-2} =\frac{-1}{2-z} =\frac{-1}2\,\frac1{1-\frac z2} =\frac{-1}2\left(1+\frac z2+\frac{z^2}4+\frac{z^3}8+\cdots\right),\qquad \left|\frac z2\right|<1\  \Longrightarrow\ |z|<2}
\end{uncoverenv}
\EMgap
\begin{uncoverenv}<4->
\EMeq{\frac1{4+z} =\frac14\,\frac1{1+\frac z4} =\frac14\left(1-\frac z4+\frac{z^2}{16}-\frac{z^3}{64}+\cdots\right),\qquad \left|-\frac z4\right|<1\  \Longrightarrow\ |z|<4}
\end{uncoverenv}
\EMgap
\begin{uncoverenv}<5->
\EMeq{\frac d{dz}\left(\frac1{4+z}\right) =\frac{-1}{(4+z)^2} =\frac14\left(-\frac14+\frac{2z}{16}-\frac{3z^2}{64}+\cdots\right)\; \Longrightarrow\;\frac1{(4+z)^2} =\frac14\left(\frac14-\frac{2z}{16}+\frac{3z^2}{64}-\cdots\right)}
\end{uncoverenv}
\EMgap
\begin{uncoverenv}<6->
\EMeq{f(z) =(4z^2+30z+68)\times\frac1{(z+4)^2}\times\frac1{z-2}}
\end{uncoverenv}
\end{EMbox}
\end{frame}

\def\EMtitlesize{24}
\begin{frame}{مثال‌هایی از بسط تیلور و لورنت (ادامه)}
\begin{EMbox}{ex}{1}{}
\begin{uncoverenv}<2->
\EMeq{f(z) =(4z^2+30z+68)\times\left(\frac{-1}2\right)\left(1+\frac z2+\frac{z^2}4+\frac{z^3}8+\cdots\right)\times\frac14\left(\frac14-\frac z8+\frac{3z^2}{64}-\cdots\right)}
\end{uncoverenv}
\EMgap
\begin{uncoverenv}<3->
\EMeq{f(z) =\frac{-1}{16}\left(34+15z+\frac{67}8z^2+\cdots\right),\qquad |z|<2}
\end{uncoverenv}
\EMgap
\begin{uncoverenv}<4->
\EMp{(ناحیهٔ همگرائی با شعاع ۲.)}
\end{uncoverenv}
\end{EMbox}
\end{frame}

\def\EMtitlesize{24}
\begin{frame}{مثال‌هایی از بسط تیلور و لورنت (ادامه)}
\begin{minipage}[t]{500bp}\vspace{0pt}
\end{minipage}
\EMside{0.0bp}{\uncover<2->{\EMfigure{1.500}{m13-radius-2}{قطب‌های \EMim{z=-4} و \EMim{z=2}؛ دایرهٔ همگرائی حول مبدأ شعاع ۲ دارد}}}
\end{frame}

\def\EMtitlesize{24}
\begin{frame}{مثال‌هایی از بسط تیلور و لورنت (ادامه)}
\begin{EMbox}{ex}{1}{}
\EMp{مطلوب است محاسبهٔ سری لورنت تابع \EMim{f(z)=\cos z^2} در همسایگی نقطهٔ \EMim{z_0=0}.}
\EMgap
\begin{uncoverenv}<2->
\EMp{با توجه به تعریف تابع کسینوس، داریم:}
\vspace{3bp}
\EMeq{\cos z =1-\frac{z^2}{2!}+\frac{z^4}{4!}-\frac{z^6}{6!}+-\cdots}
\end{uncoverenv}
\EMgap
\begin{uncoverenv}<3->
\EMeq{f(z) =\cos z^2 =1-\frac{z^4}{2!}+\frac{z^8}{4!}-\frac{z^{12}}{6!}+-\cdots}
\end{uncoverenv}
\EMgap
\begin{uncoverenv}<4->
\EMp{ناحیهٔ همگرائی تمام صفحهٔ مختلط با شعاع همگرائی بی‌نهایت.}
\end{uncoverenv}
\end{EMbox}
\end{frame}

\def\EMtitlesize{24}
\begin{frame}{مثال‌هایی از بسط تیلور و لورنت (ادامه)}
\begin{EMbox}{ex}{1}{}
\EMp{مطلوب است محاسبهٔ سری لورنت تابع \EMim{f(z)=\dfrac{\cos z}{1-z^2}} در همسایگی نقطهٔ \EMim{z_0=0}.}
\EMgap
\begin{uncoverenv}<2->
\EMeq{\left\{\begin{aligned}
\frac1{1-z^2}&=1+z^2+z^4+z^6+\cdots&&|z^2|<1\ \ \ |z|<1\\
\cos z&=1-\frac{z^2}{2!}+\frac{z^4}{4!}-\frac{z^6}{6!}+-\cdots&&|z|<\infty
\end{aligned}\right.}
\end{uncoverenv}
\EMgap
\begin{uncoverenv}<3->
\EMeq{\begin{aligned}
f(z)=\frac{\cos z}{1-z^2}&=\big(1+z^2+z^4+z^6+\cdots\big)\left(1-\frac{z^2}{2!}+\frac{z^4}{4!}-\frac{z^6}{6!}+-\cdots\right)\\
&=1+\left(1-\frac1{2!}\right)z^2+\left(1+\frac1{4!}-\frac1{2!}\right)z^4+\cdots,\qquad |z|<1
\end{aligned}}
\end{uncoverenv}
\end{EMbox}
\end{frame}

\def\EMtitlesize{24}
\begin{frame}{مثال‌هایی از بسط تیلور و لورنت (ادامه)}
\begin{EMbox}{ex}{1}{}
\EMp{مطلوب است محاسبهٔ سری لورنت تابع \EMim{f(z)=\dfrac1{z^4}} در همسایگی نقطهٔ \EMim{z_0=1}.}
\EMgap
\begin{uncoverenv}<2->
\EMeq{g(z) =\frac1z =\frac1{z-1+1} =\frac1{1+(z-1)} =1-(z-1)+(z-1)^2-(z-1)^3+-\cdots,\qquad |z-1|<1}
\end{uncoverenv}
\EMgap
\begin{uncoverenv}<3->
\EMeq{\frac d{dz}g(z) =\frac{-1}{z^2} =-1+2(z-1)-3(z-1)^2+4(z-1)^3-+\cdots,\qquad |z-1|<1}
\end{uncoverenv}
\EMgap
\begin{uncoverenv}<4->
\EMeq{\frac{d^2}{dz^2}g(z) =\frac2{z^3} =2-6(z-1)+12(z-1)^2-20(z-1)^3+-\cdots,\qquad |z-1|<1}
\end{uncoverenv}
\EMgap
\begin{uncoverenv}<5->
\EMeq{\frac{d^3}{dz^3}g(z) =\frac{-6}{z^4} =-6+24(z-1)-60(z-1)^2-+\cdots,\qquad |z-1|<1}
\end{uncoverenv}
\EMgap
\begin{uncoverenv}<6->
\EMeq{f(z) =\frac1{z^4} =1-4(z-1)+10(z-1)^2-+\cdots}
\end{uncoverenv}
\end{EMbox}
\end{frame}

\def\EMtitlesize{24}
\begin{frame}{مثال‌هایی از بسط تیلور و لورنت (ادامه)}
\begin{minipage}[t]{500bp}\vspace{0pt}
\end{minipage}
\EMside{0.0bp}{\uncover<2->{\EMfigure{1.500}{m13-radius-3}{ناحیهٔ همگرائی سری حول \EMim{z_0=1} دایره‌ای به شعاع ۱ است (تا قطب \EMim{z=0})}}}
\end{frame}

\def\EMtitlesize{24}
\begin{frame}{مثال‌هایی از بسط تیلور و لورنت (ادامه)}
\begin{EMbox}{exe}{1}{}
\EMp{مطلوب است بسط لورنت توابع زیر در همسایگی \EMim{z_0=0}:}
\EMgap
\begin{uncoverenv}<2->
\EMitem{۱.}{\dimexpr\linewidth-26bp\relax}{\EMim{f_1(z)=\dfrac1{\sqrt{1-z^2}}}}
\end{uncoverenv}
\EMgap
\begin{uncoverenv}<3->
\EMitem{۲.}{\dimexpr\linewidth-26bp\relax}{\EMim{f_2(z)=\sin^{-1}z}}
\end{uncoverenv}
\EMgap
\begin{uncoverenv}<4->
\EMitem{۳.}{\dimexpr\linewidth-26bp\relax}{\EMim{f_3(z)=\dfrac{e^z}{z^2(z^2+1)}}}
\end{uncoverenv}
\EMgap
\begin{uncoverenv}<5->
\EMitem{۴.}{\dimexpr\linewidth-26bp\relax}{تابع \EMim{\sinh z} را بر حسب توانهائی از \EMim{(z-\pi i)} بسط داده، ثابت کنید: \EMim{\displaystyle\lim_{z\to\pi i}\frac{\sinh z}{z-\pi i}=-1}}
\end{uncoverenv}
\EMgap
\begin{uncoverenv}<6->
\EMitem{۵.}{\dimexpr\linewidth-26bp\relax}{ثابت کنید: \EMim{\displaystyle z\cosh z^2=z+\sum_{n=1}^{\infty}\frac1{(2n)!}\,z^{4n+1}}}
\end{uncoverenv}
\end{EMbox}
\end{frame}

\def\EMtitlesize{24}
\begin{frame}{مثال‌هایی از نقاط تکین و مانده}
\begin{EMbox}{ex}{1}{}
\EMp{مطلوب است محاسبهٔ سری لورنت تابع \EMim{f(z)=\dfrac{z}{(z+1)(z+4)^3}} در همسایگی نقاط غیرتحلیلی آن و محاسبهٔ مانده در آن نقاط.}
\EMgap
\begin{uncoverenv}<2->
\EMp{تابع در \EMim{z=-1} و \EMim{z=-4} تحلیلی نیست.}
\end{uncoverenv}
\EMgap
\begin{uncoverenv}<3->
\EMp{\textbf{الف)} در همسایگی \EMim{z=-1} دایره‌ای به مرکز این نقطه و به شعاع کمتر از ۳ (\EMim{|z+1|<3}) نقطهٔ تکین (ایزولهٔ) \EMim{z=-4} را در بر ندارد.}
\vspace{3bp}
\EMeq{\frac{z}{(z+4)^3} =\frac{z+4-4}{(z+4)^3} =\frac{z+4}{(z+4)^3}+\frac{-4}{(z+4)^3} =\frac1{(z+4)^2}-\frac4{(z+4)^3}}
\end{uncoverenv}
\EMgap
\begin{uncoverenv}<4->
\EMeq{\frac1{z+4} =\frac1{(z+1)+3} =\frac13\,\frac1{1+\frac{z+1}3} =\frac13\left[1-\frac{z+1}3+\left(\frac{z+1}3\right)^2-\left(\frac{z+1}3\right)^3+-\cdots\right]}
\end{uncoverenv}
\EMgap
\begin{uncoverenv}<5->
\EMeq{\frac{-1}{(z+4)^2} =\frac13\left[-\frac13+\frac2{3^2}(z+1)-\frac3{3^3}(z+1)^2+-\cdots\right]\; \Longrightarrow\;\frac1{(z+4)^2} =\frac19-\frac2{27}(z+1)+\frac1{27}(z+1)^2-+\cdots}
\end{uncoverenv}
\end{EMbox}
\end{frame}

\def\EMtitlesize{24}
\begin{frame}{مثال‌هایی از نقاط تکین و مانده (ادامه)}
\begin{EMbox}{ex}{1}{}
\begin{uncoverenv}<2->
\EMeq{\frac1{(z+4)^3} =\frac16\left[\frac2{3^2}-\frac{2\times3}{3^3}(z+1)+\frac{3\times4}{3^4}(z+1)^2-\frac{4\times5}{3^5}(z+1)^3+-\cdots\right] \\[4pt] =\frac1{27}-\frac1{27}(z+1)+\frac2{81}(z+1)^2-+\cdots}
\end{uncoverenv}
\EMgap
\begin{uncoverenv}<3->
\EMeq{\frac z{(z+4)^3} =\frac1{(z+4)^2}-\frac4{(z+4)^3} =-\frac1{27}+\frac2{27}(z+1)-\frac5{81}(z+1)^2+\cdots}
\end{uncoverenv}
\EMgap
\begin{uncoverenv}<4->
\EMeq{f(z) =\frac1{z+1}\,\frac z{(z+4)^3} =-\frac1{27}\,\frac1{z+1}+\frac2{27}-\frac5{81}(z+1)+\cdots}
\end{uncoverenv}
\EMgap
\begin{uncoverenv}<5->
\EMp{نوع نقطهٔ تکین: قطب ساده (از مرتبهٔ ۱). مانده: \EMim{\operatorname{Res}\{f(z)\}\big|_{z=-1}=-\dfrac1{27}}.}
\end{uncoverenv}
\EMgap
\begin{uncoverenv}<6->
\EMp{\textbf{ب)} در همسایگی \EMim{z=-4} دایره‌ای به مرکز این نقطه و به شعاع کمتر از ۳ (\EMim{|z+4|<3}) نقطهٔ تکین (ایزولهٔ) \EMim{z=-1} را در بر ندارد.}
\vspace{3bp}
\EMeq{\frac z{z+1} =\frac{z+1-1}{z+1} =\frac{z+1}{z+1}-\frac1{z+1} =1-\frac1{z+1}}
\end{uncoverenv}
\end{EMbox}
\end{frame}

\def\EMtitlesize{24}
\begin{frame}{مثال‌هایی از نقاط تکین و مانده (ادامه)}
\begin{EMbox}{ex}{1}{}
\begin{uncoverenv}<2->
\EMeq{\frac1{z+1} =\frac1{z+4-3} =\frac1{-3+(z+4)} =\frac{-1}3\,\frac1{1-\frac{z+4}3} =\frac{-1}3\left[1+\frac{z+4}3+\frac{(z+4)^2}{3^2}+\cdots\right],\qquad |z+4|<3}
\end{uncoverenv}
\EMgap
\begin{uncoverenv}<3->
\EMeq{\frac z{z+1} =1+\frac13+\frac1{3^2}(z+4)+\frac1{3^3}(z+4)^2+\frac1{3^4}(z+4)^3+\cdots}
\end{uncoverenv}
\EMgap
\begin{uncoverenv}<4->
\EMeq{f(z) =\frac z{(z+1)(z+4)^3} =\frac1{(z+4)^3}\,\frac z{z+1} \\[4pt] =\frac{\frac43}{(z+4)^3}+\frac1{3^2}\frac1{(z+4)^2}+\frac1{3^3}\frac1{z+4}+\frac1{3^4}+\frac1{3^5}(z+4)+\frac1{3^6}(z+4)^2+\cdots}
\end{uncoverenv}
\EMgap
\begin{uncoverenv}<5->
\EMp{نوع نقطهٔ تکین: قطب از مرتبهٔ ۳. مانده: \EMim{\operatorname{Res}\{f(z)\}\big|_{z=-4}=\dfrac1{27}}.}
\end{uncoverenv}
\end{EMbox}
\end{frame}

\def\EMtitlesize{24}
\begin{frame}{مثال‌هایی از نقاط تکین و مانده (ادامه)}
\begin{minipage}[t]{500bp}\vspace{0pt}
\end{minipage}
\EMside{0.0bp}{\uncover<2->{\EMfigure{1.195}{m13-poles-1}{قطب‌های \EMim{z=-1} و \EMim{z=-4} تابع \EMim{f(z)=\frac{z}{(z+1)(z+4)^3}}}}}
\end{frame}

\def\EMtitlesize{24}
\begin{frame}{مثال‌هایی از نقاط تکین و مانده (ادامه)}
\begin{EMbox}{ex}{1}{}
\EMp{مطلوب است محاسبهٔ سری لورنت تابع \EMim{f(z)=e^{1/z}} در همسایگی نقطهٔ غیرتحلیلی آن و محاسبهٔ مانده در آن نقطه.}
\EMgap
\begin{uncoverenv}<2->
\EMeq{e^z =1+z+\frac{z^2}{2!}+\frac{z^3}{3!}+\cdots\; \Longrightarrow\;e^{1/z} =1+\frac1z+\frac1{2!\,z^2}+\frac1{3!\,z^3}+\cdots}
\end{uncoverenv}
\EMgap
\begin{uncoverenv}<3->
\EMp{نوع نقطهٔ تکین: تکین اصلی (بی‌نهایت جمله با توان‌های منفی). مانده: \EMim{\operatorname{Res}\{f(z)\}\big|_{z=0}=1}.}
\end{uncoverenv}
\end{EMbox}
\end{frame}

\def\EMtitlesize{24}
\begin{frame}{مثال‌هایی از نقاط تکین و مانده (ادامه)}
\begin{EMbox}{ex}{1}{}
\EMp{مطلوب است محاسبهٔ سری لورنت تابع \EMim{f(z)=\cos\!\left(\dfrac z{z-1}\right)} در همسایگی نقطهٔ غیرتحلیلی آن و محاسبهٔ مانده در آن نقطه.}
\EMgap
\begin{uncoverenv}<2->
\EMeq{\frac z{z-1} =1+\frac1{z-1}\; \Longrightarrow\;\cos\!\left(\frac z{z-1}\right) =\cos\!\left(1+\frac1{z-1}\right) =\cos1\cos\!\left(\frac1{z-1}\right)-\sin1\sin\!\left(\frac1{z-1}\right)}
\end{uncoverenv}
\EMgap
\begin{uncoverenv}<3->
\EMeq{\cos z =1-\frac{z^2}{2!}+\frac{z^4}{4!}-\frac{z^6}{6!}+-\cdots,\qquad \sin z =z-\frac{z^3}{3!}+\frac{z^5}{5!}-+\cdots}
\end{uncoverenv}
\EMgap
\begin{uncoverenv}<4->
\EMeq{\begin{aligned}
\cos\!\left(\frac z{z-1}\right)&=\cos1\left(1-\frac1{2!}\frac1{(z-1)^2}+\frac1{4!}\frac1{(z-1)^4}-+\cdots\right)\\
&\quad-\sin1\left(\frac1{z-1}-\frac1{3!}\frac1{(z-1)^3}+\frac1{5!}\frac1{(z-1)^5}-+\cdots\right)\\
&=\cos1-\sin1\,\frac1{z-1}-\cos1\,\frac1{2!}\frac1{(z-1)^2}+\sin1\,\frac1{3!}\frac1{(z-1)^3}+\cdots
\end{aligned}}
\end{uncoverenv}
\EMgap
\begin{uncoverenv}<5->
\EMp{نوع نقطهٔ تکین: تکین اصلی. مانده: \EMim{\operatorname{Res}\{f(z)\}\big|_{z=1}=-\sin1}.}
\end{uncoverenv}
\end{EMbox}
\end{frame}

\def\EMtitlesize{24}
\begin{frame}{مثال‌هایی از نقاط تکین و مانده (ادامه)}
\begin{EMbox}{thm}{1}{}
\begin{uncoverenv}<2->
\EMp{ثابت کنید هرگاه تابع \EMim{f(z)=\dfrac{q(z)}{g(z)}} بوده و \EMim{g(z)} در نقطهٔ \EMim{z_0} ریشهٔ مرتبهٔ اول داشته باشد، آنگاه}
\vspace{3bp}
\EMeq{\operatorname{Res}\{f(z)\}\Big|_{z=z_0} =\frac{q(z_0)}{g'(z_0)}}
\end{uncoverenv}
\end{EMbox}
\begin{EMbox}{prf}{2}{}
\begin{uncoverenv}<3->
\EMp{چون نقطهٔ \EMim{z_0} قطب مرتبهٔ ۱ است داریم (\EMim{p=1}):}
\vspace{3bp}
\EMeq{\begin{aligned}
\operatorname{Res}\{f(z)\}\Big|_{z=z_0}=A_1&=\frac1{(p-1)!}\left.\frac{d^{\,p-1}}{dz^{\,p-1}}\Big[(z-z_0)^pf(z)\Big]\right|_{z=z_0}=(z-z_0)\frac{q(z)}{g(z)}\Bigg|_{z=z_0}\\
&=(z-z_0)\frac{q(z)}{g(z)-g(z_0)}\Bigg|_{z=z_0}=\frac{q(z)}{\dfrac{g(z)-g(z_0)}{z-z_0}}\Bigg|_{z=z_0}=\frac{q(z_0)}{g'(z_0)}
\end{aligned}}
\end{uncoverenv}
\end{EMbox}
\end{frame}

\def\EMtitlesize{24}
\begin{frame}{مثال‌هایی از نقاط تکین و مانده (ادامه)}
\begin{EMbox}{exe}{1}{}
\EMp{سری لورنت توابع زیر را حول \EMim{z=0} به دست آورده، نوع این نقطه را تعیین کنید.}
\vspace{3bp}
\EMeq{f_1(z) =\frac z{e^z-1},\qquad f_2(z) =\frac{e^{-z}}{z^3},\qquad f_3(z) =\frac1{z^2(1-z^2)},\qquad f_4(z) =\frac{\sin z^2}{z^3}}
\end{EMbox}
\end{frame}
